\chapter{Potts fields, parafermions, and positive stresses}
\label{ch:factors}

The six currents explain only part of the energy in $U$. To control
the rest we need familiar rational algebras embedded in the actual
lattice realization. This chapter finds a three-state Potts factor,
a level-thirteen parafermion factor, and two useful finite extensions.
Their most important output is a positive linear expression for the
ambient conformal vector. The next chapter will turn that expression
into rationality. Here we work out the small calculations that make
the expression tangible.

\section{A parafermion is an affine algebra with its Cartan removed}

The affine algebra $A_{1,k}$ contains a rank-one Heisenberg algebra
generated by its Cartan current. Its commutant is the parafermion VOA
$K(\mathfrak{sl}_2,k)$. Subtracting the Heisenberg stress gives
\begin{equation}\label{eq:c-pf-central}
 c_{\mathrm{PF}_k}=\frac{3k}{k+2}-1
                  =\frac{2(k-1)}{k+2}.
\end{equation}
The word ``parafermion'' does not mean that this vacuum algebra contains
fractional-weight local generators. Its modules have fractional weights;
the vacuum VOA is integrally graded.

Write $L(c,h)$ for the simple highest-weight Virasoro module of central
charge $c$ and lowest weight $h$.
For $k=3$ the charge is $4/5$. The algebra is the chiral three-state
Potts theory,
\[
 P=L(4/5,0)\oplus L(4/5,3).
\]
Its extra weight-three field is essential. Finding only a Virasoro
vector of charge $4/5$ does not establish that the Potts extension
belongs to the generated algebra.

For $k=13$ the charge is $8/5$. It has a convenient presentation with
generators of weights $2,3,4,5$. The weight-four and weight-five fields
can be obtained recursively from the weight-three field, but their
singular products and decoupling relations are still needed to
recognize the algebra. A matching central charge and a single cubic
coupling would be insufficient.

\subsection{Affine fusion and the six Potts modules}
\label{sec:c-affine-pf-labels}
\index{affine fusion!level truncation}
Recall that the integrable $SU(2)_k$ modules have highest-weight
labels $\lambda=0,1,\ldots,k$, internal top spin $\lambda/2$,
and conformal weight
$h_k(\lambda)=\lambda(\lambda+2)/(4(k+2))$.
Their fusion is the level-truncated angular-momentum rule
\begin{equation}\label{eq:c-affine-truncated-fusion}
 [\lambda]\otimes[\mu]
 =\bigoplus_{\substack{\nu=|\lambda-\mu|\\
                       \nu\equiv\lambda+\mu\ (2)}}
              ^{\min(\lambda+\mu,\,2k-\lambda-\mu)}[\nu].
\end{equation}
The summation proceeds in steps of two. This is fusion of whole
affine modules, not the ordinary tensor product of their finite
current-algebra tops. For instance, at level three,
$[2]\otimes[2]=[0]\oplus[2]$: the spin-two top channel that would
occur in an ordinary spin-one tensor product is removed. The
two labels $[0],[2]$ already have the Fibonacci fusion rule.

\index{parafermion algebra!module labels}
For $\mathrm{PF}_k$, keep both the affine label and a Cartan charge:
write its modules as $\mathsf P_{\lambda,m}$, where
\begin{equation}\label{eq:c-pf-label-identification}
\begin{gathered}
 0\le\lambda\le k,\qquad m\in\Z/(2k),\qquad
 m\equiv\lambda\pmod2,\\
 (\lambda,m)\sim(k-\lambda,m+k).
\end{gathered}
\end{equation}
The parity condition follows because current modes change Cartan
charge by even integers. Charge periodicity and the last field
identification are part of the standard affine-to-parafermion
branching theorem. Equivalently, decompose the affine module under
the Cartan lattice with generator $\gamma^2=2k$:
\[
 L_k(\lambda)=
 \bigoplus_{\substack{m\bmod2k\\m\equiv\lambda\ (2)}}
 V_{\Z\gamma+m\gamma/(2k)}\otimes\mathsf P_{\lambda,m}.
\]
The charge-$m$ Heisenberg ground energy is $m^2/(4k)$.
Consequently, if its first Cartan-highest vector occurs at affine
grade $N$, its parafermion energy is
\begin{equation}\label{eq:c-pf-energy-subtraction}
 h_{\mathrm{PF}}=\frac{\lambda(\lambda+2)}{4(k+2)}
                         +N-\frac{m^2}{4k}.
\end{equation}
Dropping $N$ gives a twist modulo one, not always a lowest weight.
Fusion uses the range in \eqref{eq:c-affine-truncated-fusion} and
adds the charge labels:
\[
 \mathsf P_{\lambda,m}\otimes\mathsf P_{\mu,n}
   =\bigoplus_{\nu}\mathsf P_{\nu,m+n}.
\]
Their quantum dimensions depend only on the affine label:
\[
 d_{\lambda,m}
 =\frac{\sin((\lambda+1)\pi/(k+2))}{\sin(\pi/(k+2))}.
\]

\begin{example}[All six Potts sectors and their actual top weights]
\label{ex:c-potts-six-modules}
At $k=3$, the identification in
\eqref{eq:c-pf-label-identification} lets us use the even affine
labels $0,2$ and charges $0,2,-2$. Put
$\psi=\mathsf P_{0,2}$ and $\tau=\mathsf P_{2,0}$.
The complete table is
\begin{center}
\begin{tabular}{c|c|c|c|c}
Module & $(\lambda,m)$ & first grade $N$ & lowest weight & dimension\\
\hline
$\one$ & $(0,0)$ & $0$ & $0$ & $1$\\
$\psi$ & $(0,2)$ & $1$ & $2/3$ & $1$\\
$\psi^2$ & $(0,-2)$ & $1$ & $2/3$ & $1$\\
$\tau$ & $(2,0)$ & $0$ & $2/5$ & $\varphi$\\
$\tau\psi$ & $(2,2)$ & $0$ & $1/15$ & $\varphi$\\
$\tau\psi^2$ & $(2,-2)$ & $0$ & $1/15$ & $\varphi$
\end{tabular}
\end{center}
Here $\varphi=(1+\sqrt5)/2$. For the charged vacuum branches,
$E_{(-1)}\one$ and $F_{(-1)}\one$ are the first states of
charges $\pm2$; they have affine grade one and Heisenberg energy
$4/12=1/3$. Their parafermion weight is therefore $2/3$.
For affine label two, all three charges occur already in the
spin-one top, at affine energy $2/5$. Subtraction gives $2/5$
at charge zero and $2/5-1/3=1/15$ at charges $\pm2$.

The fusion formula gives
$\psi^3=\one$, $\tau^2=\one\oplus\tau$, and the other products
by multiplying these relations. The neutral Fibonacci subcategory
and the pointed subcategory generated by $\psi$ centralize one
another: for example their double-braiding scalar is
\[
 \frac{\theta_{\tau\psi}}{\theta_\tau\theta_\psi}
 =e^{2\pi i(1/15-2/5-2/3)}=1.
\]
The same holds for $\psi^2$. These two braided subcategories have
trivial intersection and their products exhaust the six simples,
giving
\[
 \Rep(P)\simeq\mathrm{Fib}\boxtimes
                  \mathcal C(\Z/3,-r^2/3).
\]
Here $\boxtimes$ between categories denotes their external tensor
product: its simple objects are pairs $X\boxtimes Y$, and fusion
and braiding act separately in the two factors. Dimensions and twists
therefore multiply. This differs from $M\boxtimes_V N$, which fuses
two modules inside the single category $\Rep(V)$.
For a finite abelian group $G$, the notation $\mathcal C(G,q)$ denotes
the pointed braided category with simples $g\in G$, fusion given by
group addition, and twist
$e^{2\pi i q(g)}$. The displayed quadratic form means that $\psi^r$ has twist
$e^{-2\pi i r^2/3}$; the Fibonacci object has twist $e^{4\pi i/5}$.

This table is the origin of the functions $p_1$ and $p_\tau$
used in Chapter~\ref{ch:branching}: $p_\epsilon(m)$ is the
lowest weight in the sector $\epsilon\psi^r$ with $r=m\bmod3$.
Thus a charge divisible by three gives $0$ or $2/5$, while either
nonzero residue gives $2/3$ or $1/15$. This charge-orbit label
is $r$; the Cartan label of the parafermion notation above is $2r$.
\end{example}

\subsection{What the seven even level-thirteen labels mean}
\label{sec:c-pf-grade}
\index{parafermion algebra!lowest grades}
For odd level $k$, the even affine labels form a fusion subcategory,
denoted $SO(3)_k$ here. At $k=13$ they are $2a$, $0\le a\le6$.
Equation~\eqref{eq:c-affine-truncated-fusion} becomes
\[
 a\otimes b=
 \bigoplus_{c=|a-b|}^{\min(a+b,\,13-a-b)}c,
 \qquad h_a=\frac{a(a+1)}{15}.
\]
For example, $6\otimes6=0\oplus1$. This is why the symbol
$SO(3)_{13}$ later contributes seven simple labels to the core,
and hence fourteen after combining it with Fibonacci.

The oddness of thirteen also lets us represent every
parafermion class uniquely by $(\lambda,m)=(2a,2r)$ with
$r\in\Z/13$. For the representative $-6\le r\le6$, its first
affine grade is $\max(0,|r|-a)$. Charges with $|r|\le a$ are
already in the finite spin-$a$ top. At affine grade $N$, current
descendants have finite spin at most $a+N$: each negative current
mode has spin one and costs at least one grade. Hence
$|r|\le a+N$; zero modes cannot evade this bound because they stay
inside those finite spin representations. Applying $E_{(-1)}$ to
the highest top, or $F_{(-1)}$ to the lowest top, attains the bound
in this range; affine
integrability allows these powers because
$|r|-a\le6-a\le13-2a$. The resulting vectors are killed by
positive Cartan modes, as their commutators reduce to raising or
lowering modes that kill the appropriate extreme top. Therefore
\[
 h_{\mathrm{PF}}(a,r)=\frac{a(a+1)}{15}
                   -\frac{r^2}{13}+\max(0,|r|-a).
\]
In particular the simple-current charge $a=0,r=1$ has actual
weight $1-1/13=12/13$, not $-1/13$. This same integer-grade
correction will matter when the fourteen core tops are minimized.
The thirteen invertible labels $(0,2r)$ have quadratic form
$-r^2/13$ modulo one. These are the pointed parafermion charges
used in the later order-thirteen extension.

\section{\texorpdfstring{The Potts stress in a scaled $A_2$ plane}{The Potts stress in a scaled A2 plane}}
\index{Potts algebra!embedded stress}

With the three momenta $a,b,c$ of \eqref{eq:c-u}, put
$\beta_1=b-a$, $\beta_2=c-b$, and $\beta_3=\beta_1+\beta_2$.
The Gram matrix of $\beta_1,\beta_2$ is
\[
 \begin{pmatrix}4&-2\\-2&4\end{pmatrix}.
\]
Thus they span $\sqrt2 A_2$. To keep scalar notation separate from
the Haagerup dimension, write $r=\sqrt{13}$ and $t=\sqrt{-3}$ in
the next formulas. Define
\begin{equation}\label{eq:c-potts-vector}
 p=\frac{\beta_1(-1)^2+\beta_2(-1)^2+
                      \beta_1(-1)\beta_2(-1)}{15}\one
       +\sum_{\delta=\pm\beta_1,\pm\beta_2,\pm\beta_3}
                                    p_\delta e^\delta,
\end{equation}
where
\[
\begin{array}{c|cc}
 \delta&p_\delta&p_{-\delta}\\\hline
 \beta_1&-(r+3t)/100&(-r+3t)/10\\
 \beta_2&(7+3rt)/100&(7-3rt)/100\\
 \beta_3&(-r+3t)/100&-(r+3t)/10
\end{array}
\]
The unequal opposite coefficients reflect the transported torus gauge.
They do not indicate failure of positivity.

\begin{example}[Computing the Potts central term by hand]
The inverse of the displayed Gram matrix shows that its plane stress
is
\[
 \omega_{\rm plane}=\frac16
 \bigl(\beta_1(-1)^2+\beta_2(-1)^2+
                              \beta_1(-1)\beta_2(-1)\bigr)\one.
\]
The Heisenberg part of $p$ is therefore $(2/5)\omega_{\rm plane}$.
Since the plane has charge two, this part contributes $4/25$ to
$p_{(3)}p$. Opposite momenta of squared norm four have cocycle sign $+1$.
Moreover
\[
 p_{\beta_1}p_{-\beta_1}
 =\frac{r^2-9t^2}{1000}=\frac1{25},
 \qquad
 p_{\beta_2}p_{-\beta_2}
 =\frac{49-9(rt)^2}{10000}=\frac1{25},
\]
and the third pair gives the same result. Both orders of each pair
occur, giving $6/25$. Thus
\[
 p_{(3)}p=\left(\frac4{25}+\frac6{25}\right)\one
          =\frac25\one,
\]
which is the central term $c/2$ for $c=4/5$.
\end{example}

There is an equally short reality check. The torus PCT factors for
$\beta_1,\beta_2,\beta_3$ are $10,1,10$. Hence, for instance,
$10\overline{p_{\beta_1}}=p_{-\beta_1}$, and the other pairs work
in the same way. The Heisenberg quadratic is fixed by PCT.

To obtain the full Potts algebra set
\[
 A=\beta_1-\beta_2,\quad C=\beta_1+2\beta_2,
 \quad E=-2\beta_1-\beta_2,\quad X_\delta=5p_\delta e^\delta,
\]
and use the weight-three field
\begin{align}\label{eq:c-potts-cubic}
 w_P={}&:h_Ah_Ch_E:
 -9\bigl(:h_A(X_{\beta_3}+X_{-\beta_3}):\notag\\
 &\hspace{25mm}+:h_C(X_{\beta_1}+X_{-\beta_1}):
              +:h_E(X_{\beta_2}+X_{-\beta_2}):\bigr),
\end{align}
where $h_A=A(-1)\one$, and similarly for the other Heisenberg
currents.

\begin{constructioninput}[The actual Potts embedding]
The vectors in \eqref{eq:c-potts-vector} and
\eqref{eq:c-potts-cubic} belong to $U$, are stable under its PCT
operation, and satisfy
\[
 p_{(1)}p=2p,\quad p_{(2)}p=0,\quad p_{(0)}p=Dp,
 \qquad p_{(1)}w_P=3w_P,\quad p_{(n)}w_P=0\ (n\geq2).
\]
Together with the displayed standard scaled-$A_2$ realization, these
fields generate the regular unitary Potts algebra $P$. Their intrinsic
translation agrees with ambient translation on this subalgebra.
\end{constructioninput}

The hand checks of central term and PCT explain two parts of this
input. They do not check every OPE coefficient or prove membership of
$w_P$ in the original generated algebra. Membership means an actual
linear identity in mode words of the original generators, not merely
that $w_P$ occurs in the ambient lattice theory.

\section{Recognizing the level-thirteen factor}
\index{parafermion algebra!level thirteen}

Let $z$ be the second embedded parafermion stress and let $a_3$ be its
normalized weight-three primary. We use the subscript to distinguish
this field from the lattice momentum $a$. Put
\[
 L=z,\quad b_4=(a_3)_{(1)}a_3,\quad
 d_5=(a_3)_{(1)}b_4,\quad c=\frac85,\quad\lambda=\frac{14}{473}.
\]
The self-OPE is especially instructive:
\begin{equation}\label{eq:c-pf-self}
\begin{array}{c|cccccc}
 n&5&4&3&2&1&0\\\hline
 (a_3)_{(n)}a_3&\frac8{15}\one&0&2L&DL&b_4&
                       \frac12Db_4-\frac1{12}D^3L.
\end{array}
\end{equation}
The normalizations $c/3$ and $2L$ are standard for a spin-three
primary. The last entry is constrained by skew symmetry, rather than
an independent free parameter.

The next row contains the nonlinear information:
\begin{align}
 (a_3)_{(n)}b_4&=0\quad(n\geq4),\notag\\
 (a_3)_{(3)}b_4&=\frac{69283}{2365}a_3,\notag\\
 (a_3)_{(2)}b_4&=\frac83\bigl(5-2\lambda(2+c)\bigr)Da_3,
 \label{eq:c-pf-second}\\
 (a_3)_{(1)}b_4&=d_5,\notag\\
 (a_3)_{(0)}b_4&=\frac25Dd_5+\frac{32\lambda}{5}:LDa_3:
 -\frac{48\lambda}{5}:(DL)a_3:
 +\frac2{15}\bigl(-5+2\lambda(c-1)\bigr)D^3a_3.\notag
\end{align}
The weight-six product $R=(a_3)_{(1)}d_5$ is not a new independent
generator. It obeys the decoupling relation
\begin{align}\label{eq:c-pf-decoupling}
 R={}&-\frac{148744}{473}:Lb_4:
 +\frac{566160}{473}:a_3a_3:
 -\frac{4800}{121}:L(:LL:):\notag\\
 &+\frac{3860}{11}:(D^2L)L:
 -\frac{33320}{121}:(DL)(DL):
 -\frac{489817}{4730}D^2b_4
 +\frac{2193061}{52030}D^4L.
\end{align}
This is where a finite presentation prevents an endless sequence of
new primary generators.

For completeness, the remaining singular row is fixed by
\begin{align*}
 (a_3)_{(7)}d_5&=c\bigl(55-16\lambda(2+c)\bigr)\one,
 \qquad (a_3)_{(6)}d_5=0,\\
 (a_3)_{(5)}d_5&=\frac43\bigl(175-64\lambda(2+c)\bigr)L,
 \\
 (a_3)_{(4)}d_5&=\bigl(110-32\lambda(2+c)\bigr)DL,\\
 (a_3)_{(3)}d_5&=\bigl(95-16\lambda(11+4c)\bigr)b_4
             +128\lambda:LL:\\
 &\quad+\bigl(10+8\lambda(c-1)\bigr)D^2L,\\
 (a_3)_{(2)}d_5&=64\lambda:(DL)L:
       +\bigl(75/2-4\lambda(13+5c)\bigr)Db_4\\
 &\quad+\frac{-25+8\lambda(9+5c)}{12}D^3L,\\
 (a_3)_{(0)}d_5&=\frac13DR+\frac{448}{1419}:LDb_4:
       -\frac{896}{1419}:(DL)b_4:\\
 &\quad-\frac{224}{1419}:(D^3L)L:
       -\frac{34019}{28380}D^3b_4
                      +\frac{9473}{170280}D^5L.
\end{align*}
The weight-eight, seven, and six formulas have deliberately retained
their normal ordering. Treating normal products as ordinary commutative
polynomials can change the derivative terms.

\begin{constructioninput}[Parafermion field identities]
Actual PCT-stable fields $z,a_3\in U$ satisfy the stress and primary
relations at charge $8/5$, with intrinsic translation $D$, and all
singular rows \eqref{eq:c-pf-self}--\eqref{eq:c-pf-decoupling} and
the subsequent displayed row hold as full lattice-vector equalities.
These are the corresponding defining rows of the simple
$K(\mathfrak{sl}_2,13)$ presentation.
\end{constructioninput}

Why do these rows identify the factor? An iterate identity gives
\[
 (b_4)_{(n)}v=(a_3)_{(1)}((a_3)_{(n)}v)
 -(a_3)_{(0)}((a_3)_{(n+1)}v)
 +(a_3)_{(n+1)}((a_3)_{(0)}v)
 -(a_3)_{(n)}((a_3)_{(1)}v).
\]
Taking $v=b_4,d_5$ will recover the $b_4,b_4$ and $b_4,d_5$ rows,
but that statement alone does not rule out circularity: the inner
products can be normal polynomials, whose evaluation must be available
first. We now give an order in which every required operation has
already been justified.

For this explanation abbreviate $a_3,b_4,d_5$ by $a,b,d$.
Commuting Virasoro modes through their definitions gives
\[
 L_m b=(2m+1)a_{(m+1)}a,\qquad
 L_m d=(2m+1)a_{(m+1)}b+a_{(1)}L_m b
 \quad(m\geq1).
\]
The supplied rows make these expressions linear in named fields and
derivatives. Thus Virasoro creation operators can be moved outside a
normal word using known commutators, and sesquilinearity removes
derivatives. Three non-Virasoro factors would already have weight at
least nine. The possible two-factor cores are therefore only
\[
\begin{array}{c|l}
 \text{weight bound}&\text{cores besides individual fields}\\\hline
 6&:aa:\\
 7&:aa:,\ :ab:\\
 8&:aa:,\ :ab:,\ :bb:,\ :ad:
\end{array}
\]
Changing the order of a pair uses skew symmetry and its known bracket.
No assumption that normal products commute is made.

Let $A_r$ mean that every nonnegative mode of $a$ can be evaluated
on every normal polynomial of weight at most $r$, and let $B_r$
mean the same assertion for $b$. The dependency order is
\[
\begin{aligned}
 A_6&\longrightarrow(b,b)\longrightarrow A_7\longrightarrow(b,d)\\
    &\longrightarrow B_6\longrightarrow A_8
                 \longrightarrow B_7\longrightarrow(d,d).
\end{aligned}
\]
To verify the first arrow, the only new core in $A_6$ is $:aa:$.
Its Wick integral brackets a coefficient of $[a_\lambda a]$
with $a$. That coefficient is linear in the vacuum and derivatives
of $L,b$, whose brackets with $a$ are supplied. The displayed
iterate then computes $(b,b)$ because every inner $a_{(j)}b$
has weight at most six. The new core $:ab:$ in $A_7$ needs only
this newly computed $(b,b)$ row. The same iterate with $v=d$
then computes $(b,d)$, since the inner $a$ products have weight
at most seven.

For $B_6$, the core $:aa:$ reduces by skew symmetry to $A_6$
on the coefficients of $[b_\lambda a]$; single-field inputs
now use the known $(b,b)$ and $(b,d)$ rows. In $A_8$, the new
core $:bb:$ uses $B_6$, and $:ad:$ uses the known brackets
of $L,b$ with $d$. Finally the new $:ab:$ core for $B_7$
uses $B_6$ on coefficients of $[b_\lambda a]$.

The final iterate is
\[
 d_{(n)}s=a_{(1)}(b_{(n)}s)-a_{(0)}(b_{(n+1)}s)
       +b_{(n+1)}(a_{(0)}s)-b_{(n)}(a_{(1)}s).
\]
At $s=d$, its first two terms use $A_8$ and its last two use
$B_7$. Hence it computes $(d,d)$ without calling any unknown
row. All generator brackets are fixed by the printed first rows,
including their simple poles and the weight-six decoupling.

The uniqueness argument is also concrete. Given two simple graded
algebras with this presentation, generate a subalgebra of their product
by the diagonal pairs of the named fields. Common polynomial closure
makes its degree-zero space the single diagonal vacuum line. The sum
of all proper graded ideals still misses that line, so there is a
unique maximal proper graded ideal. Both coordinate projections have
that kernel. Their simple quotients are therefore isomorphic.
Neither freeness nor linear independence of normal monomials is
needed: positive weights alone force the one-dimensional degree-zero
space of the diagonal algebra.

The two targets here really satisfy the simplicity premise.
The standard parafermion factor is simple. The actual generated
subalgebra is PCT-stable with its positive inherited form, and its
intrinsic translation and grading agree with the ambient ones on its
states. It is consequently of CFT type and simple. One can see the
last assertion from an orthogonal-ideal argument: a proper graded
ideal misses the vacuum; its invariant orthogonal complement contains
the vacuum and hence every state obtained by acting on the vacuum.
That complement is the full algebra, so the proper ideal is zero.
This argument identifies the actual embedded factor after the field
identities have been supplied; it does not replace their verification.

\section{The rank-two lattice--Potts core}

Return to the lattice momenta $a,b,c$ in \eqref{eq:c-u};
the temporary field abbreviations in the preceding calculation
are no longer in use. Set $v=a+b+c$. Directly,
\[
 v=(1,1,0,1,1,-3,1,1),\qquad
 r_0=v/3=\ell/2+h/6.
\]
The vectors $\ell,h$ are orthogonal, with squared norms $2,78$, so
$r_0^2=8/3$. Let
\[
 L_0=\Z\ell+\Z v,
 \qquad G_{L_0}=\begin{pmatrix}2&3\\3&24\end{pmatrix},
 \qquad\det L_0=39.
\]
We write $L_0$ for the lattice here to distinguish it from the
parafermion stress $L$ in the preceding section. This lattice symbol
is unrelated to the Virasoro energy operator, also conventionally
written $L_0$; a lattice coset or lattice VOA always makes the meaning clear.

\begin{constructioninput}[A charged lattice witness]
The actual fields satisfy
\[
 p_{(1)}u=\frac23u,\qquad p_{(2)}u=p_{(3)}u=0,
 \qquad u_{(-5)}(u_{(-3)}u)=-\frac35e^v.
\]
The Potts plane is orthogonal to the plane of $L_0$.
\end{constructioninput}

The last identity has output weight twelve, but its verification uses
only the leading powers of three exponential products. It does not
require enumerating an entire weight-twelve lattice space. This is a
good example of why following the algebraic route of a computation can
be much cheaper and clearer than building every ambient state.

Together with the left root currents and PCT, the witness puts
$V_{L_0}$ inside $U$. The field $u$ lies in
$V_{L_0+r_0}\otimes\psi$, where $\psi$ is the order-three Potts
simple current of weight $2/3$. Its total weight is
\[
 \frac{r_0^2}{2}+\frac23=\frac43+\frac23=2.
\]
The generated algebra is the three-sector extension
\begin{equation}\label{eq:c-core-T}
 \mathcal T=(V_{L_0}\otimes P)
 \oplus(V_{L_0+r_0}\otimes\psi)
 \oplus(V_{L_0+2r_0}\otimes\psi^2).
\end{equation}
Its stress is
\[
 \omega_{\mathcal T}=\frac14\ell(-1)^2\one
                         +\frac1{156}h(-1)^2\one+p,
 \qquad c_{\mathcal T}=2+\frac45=\frac{14}{5}.
\]
The extension is finite, unitary, and by integer-weight simple currents;
the standard extension theorem makes it strongly rational, in the sense
defined in Chapter~\ref{ch:regularity}.

\begin{example}[The discriminant arithmetic]
The class $g=h/39$ generates $L_0^*/L_0\cong\Z/39$ and has
quadratic form $q(g)=g^2/2=1/39$. Moreover $r_0=26g$ modulo
$L_0$: indeed $r_0-26g=\ell/2-h/2=2\ell-v\in L_0$.
The order-three lattice class has quadratic value
$26^2/39=1/3\pmod1$, canceled by the Potts simple-current weight
$2/3$. This is the locality condition for the extension. The
orthogonal remaining order-thirteen class $3g$ has
$q(3g)=3/13$. Consequently
\[
 \Rep(\mathcal T)\simeq
 \mathrm{Fib}\boxtimes\mathcal C(\Z/13,3j^2/13).
\]
The Fibonacci factor comes from the neutral Potts module data.
Its nontrivial object has lowest weight $2/5$ and twist
$e^{4\pi i/5}$, which specifies the braiding choice used here.
\end{example}

Sharing central charge with another familiar algebra is not recognition.
For example $G_{2,1}$ also has charge $14/5$, but $\mathcal T$ has
four currents and different representation data. Central charge is one
coordinate of a theory, not its name.

\section{An order-thirteen extension and fourteen useful modules}
\index{core algebra!fourteen modules}

The factors $\mathcal T$ and $\mathrm{PF}_{13}$ commute, giving
charge $22/5$. After choosing generators for their pointed parts, the
combined pointed quadratic form is
\[
 q(j,k)=-\frac{j^2+k^2}{13},\qquad(j,k)\in(\Z/13)^2.
\]
Replacing the old order-thirteen generator by twice itself converts the
$3j^2/13$ description above into this form, since
$3(2j)^2\equiv-j^2\pmod{13}$. Quadratic forms must be compared
up to a change of generator.

The isotropic nonzero lines satisfy $k=\pm5j$, because $5^2=-1$
modulo thirteen. Actual weight-two vectors in $U$ have simultaneous
energies
\[
 h_{\mathcal T}=14/13,\qquad h_{\mathrm{PF}_{13}}=12/13,
 \qquad h_{\rm complement}=0.
\]
Their Cartan weights are $(\pm1,\pm5)$ and $(0,\pm8)$.
These witnesses select one isotropic line. The resulting extension,
denoted $A_L$, has
\begin{equation}\label{eq:c-core-category}
 \Rep(A_L)\simeq\mathrm{Fib}\boxtimes SO(3)_{13}.
\end{equation}
There are fourteen simple modules: a Fibonacci label $1$ or $\tau$,
and one of the seven even affine labels $0,2,\ldots,12$.
Here $SO(3)_{13}$ carries the braiding inherited from this even-label
subcategory of the standard $SU(2)_{13}$ category.

Here actual lowest energies distinguish labels that twists alone would
not. In a connected unitary extension, an invertible constituent has
multiplicity at most one, and supported invertibles form a subgroup.
The nonzero witnessed charge generates its full order-thirteen line.
Both isotropic lines cannot occur: their joint span contains
nonintegral-weight charges. The line therefore specifies the extension
without an unresolved fixed-point multiplicity.

\section{The positive conformal-vector identity}
\index{positive stress identity}

Let $\omega_L,\omega_R$ be the two affine stresses and let $e,f$ be
the two commuting current-commutant stresses already constructed in
$S_2$. Their charges are
\[
 1,\quad\frac{117}{41},\quad
 \frac{2400}{1271},\quad\frac{70}{31},
\]
and their sum is eight. These four vectors are mutually orthogonal
for the invariant form; their squared norms are half their charges.
They span the current-invariant part of the actual space $S_2$.

Average $p$ and $z$ over the compact current group. A group average
is the orthogonal projection onto invariant vectors, so the projection
coefficients are obtained by dividing overlaps by squared norms. The
exact calculation in this finite-dimensional space gives
\[
\begin{array}{c|cccc}
 &\omega_L&\omega_R&e&f\\\hline
 \overline p&0&2/45&9/100&39/175\\
 \overline z&0&14/195&13/20&13/175
\end{array}
\]
where each row lists the coefficients in the four-axis basis.

\begin{constructioninput}[Invariant projection data]
The four-axis invariance statement and the two projection rows above
hold for the actual PCT-stable fields. Thus they are identities in the
finite-dimensional space $S_2$, independently of whether one has
already proved $S_2=U_2$.
\end{constructioninput}

\begin{proposition}[Positive stresses cover the ambient energy]
The following identity holds in $U$:
\begin{equation}\label{eq:c-positive-stress}
 \omega_U=\omega_L+\frac{3403}{4563}\omega_R
                    +\frac{25}{6}\overline p
                    +\frac{25}{26}\overline z.
\end{equation}
\end{proposition}
\begin{proof}
Check the four coordinates. The left coordinate is one. The right
coordinate is
\[
 \frac{3403}{4563}+\frac{25}{6}\frac2{45}
                  +\frac{25}{26}\frac{14}{195}
 =\frac{3403+845+315}{4563}=1.
\]
The remaining coordinates are
\[
 \frac{25}{6}\frac9{100}+\frac{25}{26}\frac{13}{20}
 =\frac38+\frac58=1,
\]
and
\[
 \frac{25}{6}\frac{39}{175}+\frac{25}{26}\frac{13}{175}
 =\frac{13}{14}+\frac1{14}=1.
\]
The four stresses sum to $\omega_U$, proving the identity.
\end{proof}

The bars should not be discarded. An averaged stress is generally not
itself a Virasoro vector, and the many conjugate Potts and parafermion
copies need not commute with one another. What survives is positivity.
In the finite-dimensional real PCT-fixed subspace, the Haar average
belongs to the convex hull of the compact orbit. Carath\'eodory's
theorem expresses it as a finite convex combination. Therefore
\begin{equation}\label{eq:c-finite-positive}
 \omega_U=\sum_{i=1}^N t_i e_i,\qquad t_i>0,
\end{equation}
where each $e_i$ is the conformal vector of an actual regular unitary
subalgebra: one of the affine factors, or a compact conjugate of Potts
or PF$_{13}$. Zero coefficients can be omitted.

This is not a tensor decomposition. It is a positive energy identity.
The next chapter uses inequalities and traces that remain valid even
when the different partial energy operators do not commute. That is
why the averaged construction is sufficient.

\section{Exercises}

\begin{exercise}\label{ex:c-factors-1}
Compute the central charges of PF$_2$, PF$_3$, and PF$_{13}$ from
\eqref{eq:c-pf-central}. Identify which has the Ising charge.

\par\smallskip\noindent{\footnotesize\normalfont
\hyperref[sol:ex:c-factors-1]{Solution on page~\pageref*{sol:ex:c-factors-1}.}}\par
\end{exercise}
\begin{exercise}\label{ex:c-factors-2}
Verify all three opposite-root products in the Potts central-term
calculation, and check PCT reality for all three momentum pairs.

\par\smallskip\noindent{\footnotesize\normalfont
\hyperref[sol:ex:c-factors-2]{Solution on page~\pageref*{sol:ex:c-factors-2}.}}\par
\end{exercise}
\begin{exercise}\label{ex:c-factors-3}
Use skew symmetry of the weight-three self-bracket to derive the simple
pole $\tfrac12Db_4-\tfrac1{12}D^3L$ in
\eqref{eq:c-pf-self}.

\par\smallskip\noindent{\footnotesize\normalfont
\hyperref[sol:ex:c-factors-3]{Solution on page~\pageref*{sol:ex:c-factors-3}.}}\par
\end{exercise}
\begin{exercise}\label{ex:c-factors-4}
Show that the two isotropic lines of $-(j^2+k^2)/13$ are exactly
$k=\pm5j$. Explain why adjoining both lines would violate integral
vacuum grading.

\par\smallskip\noindent{\footnotesize\normalfont
\hyperref[sol:ex:c-factors-4]{Solution on page~\pageref*{sol:ex:c-factors-4}.}}\par
\end{exercise}
\begin{exercise}\label{ex:c-factors-5}
Verify all four coefficients of \eqref{eq:c-positive-stress}.
Explain why checking the norms of $p$ and $z$ alone would not prove
this identity or establish either embedded factor.

\par\smallskip\noindent{\footnotesize\normalfont
\hyperref[sol:ex:c-factors-5]{Solution on page~\pageref*{sol:ex:c-factors-5}.}}\par
\end{exercise}
\begin{exercise}\label{ex:c-factors-6}
Give a finite-dimensional example of two noncommuting positive
self-adjoint matrices. Show that their positive linear combination
still gives a valid expectation-value inequality. Which part of the
next chapter's argument should therefore avoid simultaneous
diagonalization?

\par\smallskip\noindent{\footnotesize\normalfont
\hyperref[sol:ex:c-factors-6]{Solution on page~\pageref*{sol:ex:c-factors-6}.}}\par
\end{exercise}
